\section*{Introduction}
L'objectif de ce projet est de concevoir un système de résolution de noms de domaine (DNS).
Cette implémentation doit se faire en python et à partir de zéro, sans utiliser de bibliothèques spécialisées.
Chaque serveur du système doit être capable de résoudre les requêtes DNS de manière récursive et posséder un cache en plus de gérer l'éventuelle zone dont elle serait responsable.

Ce projet à réaliser par équipe de 1 étudiant, demande d'appliquer la méthodologie de gestion de projet \textit{Agile} pour la conception et l'implémentation du système DNS.
En particulier, il est demandé d'implémenter des tests unitaires pour valider le bon fonctionnement du système et garantir la qualité du code produit.
Les artéfacts produits au cours du projet sont constatables sur le dépôt GitLab~\cite{gitlab_dns_project} et sont constitués de la documentation, du code source et des tests unitaires.

Le présent rapport décrit la conception et l'implémentation de ce système DNS.
Il commencera par une analyse du contexte et du cahier des charges, suivie d'une analyse fonctionnelle listant les besoins et les contraintes du système.
L'analyse technique détaillera ensuite l'architecture logicielle et les choix de conception.
Enfin, la section sur l'implémentation et la validation présentera les résultats des tests.